\section*{Contoh}
\addcontentsline{toc}{section}{Contoh}

\begin{contoh}[Dua pohon urai]
Untuk grammar \(E \to E + E \mid E * E \mid \code{id}\) dan string
\(\code{id} + \code{id} * \code{id}\), pohon pertama memiliki akar
\code{+} dengan anak kanan \code{*}; pohon kedua memiliki akar
\code{*} dengan anak kiri \code{+}. Keduanya valid menurut grammar,
tetapi hanya pohon pertama yang sesuai dengan aturan precedence
perkalian lebih tinggi.
\end{contoh}

\begin{contoh}[Menghapus left recursion langsung]
Aturan
\[
A \to A a \mid A b \mid c \mid d
\]
diubah menjadi
\[
A \to c A' \mid d A', \qquad
A' \to a A' \mid b A' \mid \varepsilon.
\]
String \(c a b\) diturunkan sebagai
\(A \Rightarrow cA' \Rightarrow caA' \Rightarrow cabA'
\Rightarrow cab\).
\end{contoh}

\begin{contoh}[Left recursion tidak langsung]
Diberikan
\[
A \to B a \mid c, \qquad B \to A b \mid d.
\]
Substitusi \(B\) ke aturan \(A\) menghasilkan
\[
A \to A b a \mid d a \mid c.
\]
Setelah itu, left recursion langsung dihapus:
\[
A \to d a A' \mid c A', \qquad
A' \to b a A' \mid \varepsilon.
\]
\end{contoh}

\begin{contoh}[Left factoring if-else]
Aturan
\[
S \to \text{if}\ E\ \text{then}\ S\ \text{else}\ S
\mid \text{if}\ E\ \text{then}\ S
\]
difaktorkan menjadi
\[
S \to \text{if}\ E\ \text{then}\ S\ S', \qquad
S' \to \text{else}\ S \mid \varepsilon.
\]
Parser memutuskan \code{else} setelah selesai memproses \(S\).
\end{contoh}

\begin{contoh}[Dangling else]
Pada string
\code{if (a) if (b) x; else y;}, aturan else terdekat memasangkan
\code{else y;} dengan \code{if (b)}. Pohon yang dihasilkan adalah
pohon kiri pada \figref{fig:ch09-dangling}. Jika diinginkan
pasangan luar, program harus ditulis dengan blok eksplisit:
\code{if (a) \{ if (b) x; \} else y;}.
\end{contoh}
